% ===== 原文 pp.15–17（PDF pp.16–18）=====
% [原文 p.15]
Strominger–Yau–Zaslow [9] 使用 D-膜论证了 Calabi–Yau 流形之间的镜像对称是同一种意义上的
T-对偶。也就是说，镜像对中的每一个都具有拉格朗日环面纤维化的结构，而镜像对称是沿纤维的 T-对偶，
从而镜像流形可以作为缠绕在原流形环面纤维上的 D-膜的模空间而得到。\footnote{在 [9] 中还加入了另一个
重要要素——D-膜的稳定性——并论证 Calabi–Yau 流形具有\emph{特殊}拉格朗日环面纤维化的结构。}
一般而言，Calabi–Yau 流形并不容许光滑的环面纤维化，而对 (3.9) 那样的规则将存在量子修正，这些修正
在奇异纤维附近是显著的。主要问题是如何正确地计算这些量子修正。

\subsection{一个谜题与一个解答}

作为具有奇异纤维的拉格朗日环面纤维化的最简单例子，让我们考虑靶空间 $X$ 是拓扑为二维球面的复射影
直线 $\CP^1$ 的情形。它具有线段上的圆纤维化结构，如图 1 所示。

\begin{figure}[htbp]
\centering
\includegraphics[width=0.40\textwidth]{fig1.png}
\figtitle{1}{$\CP^1$ 作为线段上的圆纤维化。}
\end{figure}

注意 $\CP^1$ 不是 Calabi–Yau 流形，因而此设置与 SYZ 并不完全相同，但我们仍可问同样的问题：当对
圆纤维施加 T-对偶时会发生什么？来自奇异纤维的量子修正又是什么？圆的大小在线段中点处最大，随着离开
中点而减小，并在两端处消失。由于 T-对偶反演圆的半径 (3.4)，在对偶一侧，圆的大小在线段中点处最小，
随着离开中点而增大，并在两端处发散。因此，对偶空间 $\tX$ 应看起来像图 2，具有圆柱的拓扑。若认真对待
这一点，那么靶空间为 $X=\CP^1$ 的非线性 sigma 模型必须对偶于靶空间为图 2 中的 $\tX$ 的非线性
sigma 模型。

% [原文 p.16]
\begin{figure}[htbp]
\centering
\includegraphics[width=0.36\textwidth]{fig2.png}
\figtitle{2}{纤维化的朴素 T-对偶。}
\end{figure}

这正确吗？作为检验，让我们比较两个理论的守恒量。原系统具有旋转圆纤维的 $U(1)$ 对称性，因而有相应
的 Noether 荷——圆方向上的动量。在对偶一侧，空间 $\tX$ 在线段两端有两个洞，因而圆方向上的绕数是
守恒的。这与 T-对偶把一个理论的动量映为对偶的绕数这一事实一致。另一方面，原系统中没有守恒的绕数，
因为 $S^2$ 是单连通的——绕圆纤维的一个圈可以通过走到线段的任一端而收缩。在对偶一侧，似乎存在一个
旋转圆纤维的 $U(1)$ 对称性，因而动量似乎是守恒的。于是，两个理论的对称性并不匹配！两个理论之间还有
另一个显著差异：$\CP^1$ 是紧的，而 $\tX$ 是非紧的。当表述在圆柱 $\Sigma=\R\times S^1$ 上时，这
意味着 $\CP^1$ sigma 模型的能谱必定是离散的，而 $\tX$ sigma 模型的谱是连续的。当两个理论等价时，
这不可能。所以，朴素的 T-对偶似乎并不成立。

事实上，这一朴素的 T-对偶存在一个重要的量子修正，它来自圆纤维的退化。这就是超势的产生。设 $y$ 是
$\tX$ 的复坐标，具有周期性 $y\equiv y+2\pi\ii$，使得 $\Im(y)$ 参数化圆纤维，而 $\Re(y)$ 参数化
水平线段。结果表明 $\Re(y)$ 可以取遍整个实轴 $\R$，其中线段的左端与右端分别对应于 $-\infty$ 与
$+\infty$。于是，对偶理论具有超势 [10]
\begin{equation}
  W=\ee^{-y}+\ee^{y}. \tag{3.23}
\end{equation}
这一效应在奇异纤维的对偶附近确实很大——第一项与第二项分别在左端与右端发散。该超势的存在解决了朴素
T-对偶的两个问题。第一，这个超势
% [原文 p.17]
在 $\Im(y)$ 的平移下并不是不变的，因而在对偶理论中纤维动量不守恒。这对应于 $\CP^1$ sigma 模型中绕数的
不守恒。第二，相应的势 $U(y)=g^{y\overline{y}}|\partial_yW|^2$ 是固有的，即对每个 $E$，子集
$\{U(y)<E\}$ 是紧的。这意味着表述在 $\Sigma=\R\times S^1$ 上的理论的谱是离散的，正如 $\CP^1$
sigma 模型那样。

这引导我们断言：$\CP^1$ sigma 模型的镜像是具有圆柱靶与超势 (3.23) 的 Landau–Ginzburg 模型。
超势的产生将在第 5 节推导。正如我们将在那里看到的，这一镜像构造可推广到更广的靶空间类，包括环面簇
以及环面簇中的超平面或其完全交。推导中的主要工具是一类称为\emph{规范线性 sigma 模型}的超对称规范
理论——接下来我们就转向它。

% [Tutor Note: p.15 图 1 与 §3.3 的论证可以定量对上：把 $\CP^1$ 看成线段上的圆纤维化，纤维半径
%  $\rho(x)$ 在中点最大、两端 $\to0$；(3.4) 说 $R\to1/R$，所以对偶纤维半径 $1/\rho(x)$ 在中点
%  最小、两端 $\to\infty$，对偶空间因而“像图 2 那样具有圆柱的拓扑”。直观上，两张图是同一个线段，
%  只是把两端的退化纤维“拉长”成了两个洞——这正是正文后文说 $\tX$ “在线段两端有两个洞”的由来。]
% [Tutor Note: p.16 式 (3.23) 为何能解决第二个问题（谱离散）：在圆柱上用复坐标 $y$（$y\equiv y+2\pi\ii$），
%  $\partial_yW=-\ee^{-y}+\ee^{y}$，故 $U=g^{y\overline{y}}|\partial_yW|^2\propto|\ee^{y}-\ee^{-y}|^2$，
%  当 $\Re(y)\to\pm\infty$ 时按 $\ee^{2|\Re y|}$ 增长，于是 $\{U(y)<E\}$ 对每个 $E$ 都是紧的（proper），
%  谱必然离散。这可用来核对原文“the associated potential ... is proper”一句不是空话。]
% [Tutor Note: p.16 的“对称性不匹配”是 SYZ 式 T-对偶在奇异纤维处的失效点，两条判据都可独立验证：
%  (i) 原系统 $\CP^1=S^2$ 单连通，绕纤维的圈可滑到线段端点收缩，故无守恒绕数；对偶侧 $\tX$ 有两个洞，
%  绕数守恒，于是原系统的动量守恒在对偶侧找不到对应；(ii) 原空间紧、对偶空间非紧，紧空间上
%  $\Sigma=\R\times S^1$ 的能谱离散而非紧空间上连续，等价理论不可能一个离散一个连续。]
% [Tutor Note: p.16 原文“the naïver T-duality”是笔误（naïve 的比较级应为“more naïve”），译文作
%  “朴素的 T-对偶”；p.13 原文“at the same positon in Re(φ)”“at a fixed poition in the circle”
%  中的 positon、poition 均为 position 的笔误。另 p.15 原文“empolyed”“Lagragian”分别为
%  employed、Lagrangian 的笔误。]
% [Tutor Note: 术语首现：“proper”按冻结表译作“固有”（p.17 的 $U(y)$ 是固有势）；p.17 末句斜体的
%  “gauged linear sigma models”译作“规范线性 sigma 模型”，正是下一节（§4）的主题，译文在此保留
%  斜体以对应原文的重点标记。]
